% GENERATED FILE. Edit canonical lessons, then run npm run generate:books.
% canonical-source-hash: fnv1a64:d3cd0efe5d15e2c5
% canonical-lessons: ML-C49-true, ML-C49-enough, ML-C49-certainly, ML-C49-perhaps, ML-C49-just-so

\chapter{Short Replies}
\label{ch:ml-replies}
\hlchaptermodality{49}

\begin{chapteropening}
\textbf{By the end of this chapter:} \emph{I can answer with five short replies in Malayalam, from a flat agreement to a hedge.}

Complete ML-C49-just-so and say all five words in order.
\end{chapteropening}

\section[satyaṁ]{\ml{സത്യം} (satyaṁ) — true, the truth}

\label{lesson:ML-C49-true}



\begin{quote}
Before the new run of words: what did \emph{karṣakan} mean, and what did \emph{atithi} mean?
\end{quote}

\subsection*{You'll want to know: \ml{സത്യം}}
\textbf{\ml{സത്യം}} (\emph{satyaṁ}) — \textquotedblleft{}true, the truth\textquotedblright{}.

Sanskrit \emph{satya}, built on \emph{sat}, 'being, that which is'. Truth, in that word's picture, is simply what is the case — a noun made out of the verb 'to be'.

Malayalam already has \ml{ശരി} (\emph{śari}) for the everyday \textquotedblleft{}right, okay\textquotedblright{} of a conversation. \ml{സത്യം} is the heavier one: not 'that will do' but 'that is so'.

Kannada answers with \ka{ನಿಜ} (\emph{nija}), a different Sanskrit word, whose older sense is 'one's own, innate'. Two languages, two borrowings, and each picked a different picture of what being true is.

\begin{grammarlens}[title={What you've built}]
The first of five short replies.
\end{grammarlens}

\subsection*{Your turn}
Say these aloud:

\begin{itemize}
\raggedright
  \item \emph{satyaṁ}
  \item it once more, slowly
  \item \emph{satyaṁ}, then \emph{śari}, and hear the weight shift between them
\end{itemize}

\begin{tcolorbox}[breakable,colback=teal!4,colframe=teal!35!black,title={Before you move on}]
What does \ml{സത്യം} mean? (\textquotedblleft{}true, the truth\textquotedblright{}.) Say it once more.
\end{tcolorbox}
\section[mati]{\ml{മതി} (mati) — enough, that will do}

\label{lesson:ML-C49-enough}



\begin{quote}
Before the new one: what did \emph{satyaṁ} mean?
\end{quote}

\subsection*{You'll want to know: \ml{മതി}}
\textbf{\ml{മതി}} (\emph{mati}) — \textquotedblleft{}enough, that will do\textquotedblright{}.

Two syllables that make a whole sentence. Hand someone rice and they say \ml{മതി} and nothing else, and the exchange is complete.

Malayalam has a small handful of these — words that are already a finished utterance and take no verb, no subject and no ending. This book has met one before in \ml{ഇല്ല}.

Dictionaries connect \ml{മതി} with a root about measuring and reckoning; the connection is plausible rather than settled, and it is more honest to leave it that way than to tidy it. Kannada says \ka{ಸಾಕು} (\emph{sāku}), which is not this word at all.

\begin{grammarlens}[title={What you've built}]
Two replies: true, and enough.
\end{grammarlens}

\subsection*{Your turn}
Say these aloud:

\begin{itemize}
\raggedright
  \item \emph{mati}
  \item it once more, slowly
  \item \emph{mati}, then \emph{illa}, two words that each finish a sentence alone
\end{itemize}

\begin{tcolorbox}[breakable,colback=teal!4,colframe=teal!35!black,title={Before you move on}]
What does \ml{മതി} mean? (\textquotedblleft{}enough, that will do\textquotedblright{}.) Say it once more.
\end{tcolorbox}
\section[tīrccayāyuṁ]{\ml{തീർച്ചയായും} (tīrccayāyuṁ) — certainly, for sure}

\label{lesson:ML-C49-certainly}



\begin{quote}
Before the new one: what did \emph{mati} mean?
\end{quote}

\subsection*{You'll want to know: \ml{തീർച്ചയായും}}
\textbf{\ml{തീർച്ചയായും}} (\emph{tīrccayāyuṁ}) — \textquotedblleft{}certainly, for sure\textquotedblright{}.

A long word, and every piece of it is Malayalam's own. \ml{തീർച്ച} (\emph{tīrcca}) is 'a settling, a thing decided', off the verb \ml{തീരുക}, 'to come to an end'. Add the adverb-making tail and the whole says 'in a settled way'.

Kannada reaches for \ka{ಖಂಡಿತ} (\emph{khaṇḍita}) here, Sanskrit for 'cut, broken off'. The two languages arrived at the same idea from the same direction — a matter is certain once it has been finished off — but one built the word at home and one imported it.

There is a chillu \ml{ർ} inside, the third chapter running to end up leaning on one.

\begin{grammarlens}[title={What you've built}]
Three. True, enough, certainly.
\end{grammarlens}

\subsection*{Your turn}
Say these aloud:

\begin{itemize}
\raggedright
  \item \emph{tīrccayāyuṁ}
  \item it once more, slowly
  \item \emph{tīrccayāyuṁ}, slowly, then \emph{satyaṁ}, and put the long word beside the short one
\end{itemize}

\begin{tcolorbox}[breakable,colback=teal!4,colframe=teal!35!black,title={Before you move on}]
What does \ml{തീർച്ചയായും} mean? (\textquotedblleft{}certainly, for sure\textquotedblright{}.) Say it once more.
\end{tcolorbox}
\section[orupakṣē]{\ml{ഒരുപക്ഷേ} (orupakṣē) — perhaps}

\label{lesson:ML-C49-perhaps}



\begin{quote}
Before the new one: what did \emph{tīrccayāyuṁ} mean?
\end{quote}

\subsection*{You'll want to know: \ml{ഒരുപക്ഷേ}}
\textbf{\ml{ഒരുപക്ഷേ}} (\emph{orupakṣē}) — \textquotedblleft{}perhaps\textquotedblright{}.

Take it apart and it says \textquotedblleft{}on one side\textquotedblright{}. \ml{ഒരു} (\emph{oru}) is 'one', the same numeral this book counted with. \ml{പക്ഷം} (\emph{pakṣaṁ}) is Sanskrit for a side, a wing, one of the alternatives on offer.

So a Malayalam speaker hedges by naming a side rather than by naming a doubt. The word is half native, half borrowed, and joined so smoothly that the seam is easy to miss.

Kannada hedges with \ka{ಬಹುಶಃ} (\emph{bahuśaḥ}), a wholly Sanskrit word meaning something closer to 'in many ways'. Same job, different picture again.

\begin{grammarlens}[title={What you've built}]
Four. True, enough, certainly, perhaps.
\end{grammarlens}

\subsection*{Your turn}
Say these aloud:

\begin{itemize}
\raggedright
  \item \emph{orupakṣē}
  \item it once more, slowly
  \item \emph{orupakṣē}, then \emph{tīrccayāyuṁ}, the hedge and the promise
\end{itemize}

\begin{tcolorbox}[breakable,colback=teal!4,colframe=teal!35!black,title={Before you move on}]
What does \ml{ഒരുപക്ഷേ} mean? (\textquotedblleft{}perhaps\textquotedblright{}.) Say it once more.
\end{tcolorbox}
\section[aṅṅane]{\ml{അങ്ങനെ} (aṅṅane) — that way, just so}

\label{lesson:ML-C49-just-so}



\begin{quote}
Before the new one: what did \emph{orupakṣē} mean?
\end{quote}

\subsection*{You'll want to know: \ml{അങ്ങനെ}}
\textbf{\ml{അങ്ങനെ}} (\emph{aṅṅane}) — \textquotedblleft{}that way, just so\textquotedblright{}.

This one the book has half-taught already. The chapter on pointing laid out a pattern: \emph{i-} for what is near, \emph{a-} for what is far, \emph{e-} for what is being asked about. \ml{എങ്ങനെ} (\emph{eṅṅane}), 'how?', was the question member of exactly this set.

Swap its opening for the far one and the question becomes an answer. \ml{അങ്ങനെ} is 'that way' — the manner just spoken of, agreed to.

Kannada does the same trick with \ka{ಹಾಗೆ} (\emph{hāge}), built on its own far-pointer. The machinery is inherited; each sister runs it with her own parts.

\begin{grammarlens}[title={What you've built}]
Five, and the run is closed: true, enough, certainly, perhaps, just so.
\end{grammarlens}

\subsection*{Your turn}
Say these aloud:

\begin{itemize}
\raggedright
  \item \emph{aṅṅane}
  \item it once more, slowly
  \item all five in order — \emph{satyaṁ}, \emph{mati}, \emph{tīrccayāyuṁ}, \emph{orupakṣē}, \emph{aṅṅane}
\end{itemize}

\begin{tcolorbox}[breakable,colback=teal!4,colframe=teal!35!black,title={Before you move on}]
What does \ml{അങ്ങനെ} mean? (\textquotedblleft{}that way, just so\textquotedblright{}.) Say it once more.
\end{tcolorbox}
